
\documentclass[12pt]{article}

\usepackage[margin=1in]{geometry}
\usepackage{amsmath,amsthm,amssymb,scrextend,mathtools}
\usepackage{fancyhdr}
\usepackage{graphicx}
\usepackage{tikz}
\usepackage[mathscr]{euscript}
\let\euscr\mathscr \let\mathscr\relax
\usepackage[scr]{rsfso}
\usepackage{multicol}
\usepackage[colorlinks=true,pdfstartview=FitV,linkcolor=blue,citecolor=blue,urlcolor=blue]{hyperref}
\pagestyle{fancy}

\usepackage{listings}
\usepackage{xcolor}
\definecolor{codegreen}{rgb}{0,0.6,0}
\definecolor{codegray}{rgb}{0.5,0.5,0.5}
\definecolor{codepurple}{rgb}{0.58,0,0.82}
\definecolor{backcolour}{rgb}{0.95,0.95,0.92}
\lstdefinestyle{mystyle}{
    commentstyle=\color{codegreen},
    keywordstyle=\color{blue},
    numberstyle=\tiny\color{codegray},
    stringstyle=\color{codepurple},
    basicstyle=\ttfamily\footnotesize,
    breakatwhitespace=false,
    breaklines=true,
    captionpos=b,
    keepspaces=true,
    numbers=left,
    numbersep=5pt,
    showspaces=false,
    showstringspaces=false,
    showtabs=false,
    tabsize=2
}
\lstset{style=mystyle}

\let\ds\displaystyle
\newcommand{\cont}{\subseteq}
\newcommand{\un}{\cup}
\newcommand{\inter}{\cap}
\newcommand{\R}{\mathcal{R}}
\newcommand{\N}{\mathcal{N}}
\newcommand{\E}{\mathbb{E}}

\begin{document}

\lhead{PhD Macro: Replication Exercise}
\chead{Logan Casey}
\rhead{\today}

\tableofcontents
\newpage

\section{Paper Overview}
\subsection{What Does This Paper Do?}
\begin{itemize}
  \item Auclert-Rognlie-Straub style heterogeneous-agent model in an open-economy setting.
  \item Contractionary depreciations can arise when the real-income channel (import prices up, output down) dominates the Keynesian multiplier channel (foreign demand up, output up).
  \item These heterogeneous-agent channels can jointly dominate the representative-agent complete-markets expenditure-switching channel.
  \item The paper delivers a calibrated quantitative model and policy counterfactuals.
\end{itemize}

\section{Base Model}
\subsection{Core Ingredients}
\begin{itemize}
  \item Domestic households face idiosyncratic income shocks and a borrowing constraint.
  \item Households can hold a domestic bond, foreign bond, and domestic equity.
  \item Wages are sticky.
  \item Impulse responses are computed to MIT shocks under perfect foresight.
\end{itemize}

\subsection{Simplifications}
\begin{itemize}
  \item No illiquid asset or adjustment friction, so UIP holds and portfolio choice is indeterminate.
  \item No need to model consumption bundles in the base setup.
  \item Rest of world is representative-agent complete-markets with flexible prices and wages.
  \item No capital, taxes, debt, or government spending in the base setup.
  \item Producer-currency pricing in the baseline (relaxed in the appendix of the slide deck).
\end{itemize}

\subsection{Solution Method (SSJ)}
\begin{itemize}
  \item Solve for a steady state where all residuals are zero.
  \item Compute sequence-space Jacobians, including policy and distribution responses in the household block.
  \item Solve for unknown aggregate paths that satisfy equilibrium targets.
\end{itemize}

\noindent Unknowns:
\[
\{Y_t,\;W_t,\;Q_t,\;E_t,\;i_t,\;r_t\}
\]

\noindent Targets:
\[
\text{goods market clearing,\ NKPC,\ UIP,\ exchange rate identity,\ Fisher equation,\ MP rule}
\]

\subsection{Base Model DAG}
\begin{figure}[h!]
  \centering
  \includegraphics[width=\linewidth,height=0.75\textheight,keepaspectratio]{../figures/ha_oe_base_dag.png}
  \caption{Base model DAG}
\end{figure}

\clearpage
\section{Quantitative Extension}
\begin{itemize}
  \item Delayed substitution: consumers adjust consumption baskets with Calvo-style probability.
  \item Heterogeneous labor-shock incidence, with larger impacts on poorer households.
  \item Non-homothetic (Stone-Geary) consumption where poorer households consume more imports.
  \item Sticky prices in the quantitative extension.
\end{itemize}

\subsection{Figure 8: Devaluation Shock IRFs}
\begin{figure}[h!]
  \centering
  \includegraphics[width=\linewidth,height=0.72\textheight,keepaspectratio]{../figures/figure8_original.png}
  \caption{Paper Figure 8 (original)}
\end{figure}

\begin{figure}[h!]
  \centering
  \includegraphics[width=\linewidth,height=0.72\textheight,keepaspectratio]{../figures/ha_oe_quant_irf_policy_compare.png}
  \caption{Replication Figure 8}
\end{figure}

\clearpage
\subsection{Figure 9: Monetary Policy Tradeoffs}
\begin{figure}[h!]
  \centering
  \includegraphics[width=\linewidth,height=0.72\textheight,keepaspectratio]{../figures/figure9_original.png}
  \caption{Paper Figure 9 (original)}
\end{figure}

\begin{figure}[h!]
  \centering
  \includegraphics[width=\linewidth,height=0.72\textheight,keepaspectratio]{../figures/ha_oe_quant_policy_rqy.png}
  \caption{Replication Figure 9}
\end{figure}

\clearpage
\section{Equations and Block-Map Notes}
\subsection{Equilibrium Objects}
Given sequences of rest-of-world discount shocks $\{B_t\}$, monetary shocks $\{\varepsilon_t\}$, and initial wealth distribution $D^i_0(s,B_H,B_F,e)$, equilibrium is a path of:
\begin{itemize}
  \item Policies $\{c^H_{i,t}(a',e), c^F_{i,t}(a',e), c_{i,t}(a',e), a_{i,t+1}(a',e)\}$.
  \item Distributions $D_{i,t}(a',e)$.
  \item Prices $\{E_t,Q_t,P_t,P^H_t,P^F_t,W_t,r_t,i_t\}$.
  \item Aggregate quantities $\{C_t,C^H_t,C^F_t,Y_t,A_t,p_t,d_t,nfa_t\}$.
\end{itemize}
such that households optimize, distributions evolve consistently, firms optimize, and domestic goods markets clear.

\subsection{Steady State}
\begin{itemize}
  \item Zero net foreign assets.
  \item Zero inflation.
  \item Normalize prices, output, and consumption to one.
\end{itemize}

\subsection{Parameters}
\begin{itemize}
  \item $\beta$: discount factor.
  \item $\sigma$: inverse EIS.
  \item $\phi_{labor}$: inverse Frisch elasticity.
  \item $\eta$: home/foreign goods elasticity.
  \item $\alpha$: openness/home bias.
  \item $\psi$: labor-supply scaling parameter.
  \item $\kappa_w$: wage Phillips slope, depends on $\beta$, $\theta_w$ (Calvo probability), and $\mu_w$ (wage markdown).
  \item Not calibrated directly: $\chi=\eta(1-\alpha)+\gamma$; with $\gamma=\eta$, $\chi=\eta(2-\alpha)$.
  \item $\phi_{\pi}$: Taylor-rule inflation coefficient.
  \item $\rho_m$: monetary-shock persistence.
  \item $\underline{c}$: subsistence consumption level.
  \item $\zeta$: labor-incidence parameter in the quantitative extension.
  \item $\mu$: markup.
  \item $\theta_H$, $\kappa_H$: home-goods Calvo parameter and Phillips slope.
  \item Foreign-sector $\theta$ and $\kappa$ analogs for foreign goods and home exports.
  \item $\theta_{sub}$: consumption-reallocation (delayed-substitution) parameter.
\end{itemize}

\subsection{Base Equations (1)--(23): Mapping Notes}
\begin{itemize}
  \item (1), (2): household value and utility setup (handled inside the household block).
  \item (3): CES aggregator replaced by Stone-Geary form in quantitative version.
  \item (4): price index block (\texttt{price\_index}).
  \item (5): nominal return equalization not used directly in this simplified map.
  \item (6): real exchange-rate identity (\texttt{exchange\_rate\_identity} target residual).
  \item Fisher equation uses an adjusted form to match constant-$r$ policy setup.
  \item (7): real UIP (\texttt{uip\_real} target residual).
  \item (8): real asset-position bookkeeping is implicit in household aggregation.
  \item (9): one-asset household Bellman system (implemented in SSJ household routines).
  \item (10)--(11): policy mapping from total consumption to home/foreign components not solved as separate aggregate unknowns in the simplified map.
  \item (12): foreign demand block simplified when solving at aggregate H/F level.
  \item (13): production $Y_t=N_t$ (\texttt{production}).
  \item (14): static markup pricing omitted when Phillips blocks are used instead.
  \item (15): dividends/profits useful for IRF accounting.
  \item (16): foreign monetary block with discount-shock variation in rest-of-world rates.
  \item Import price pass-through uses $P^F_t=E_t$ in the PCP baseline.
  \item (17): loop for export pricing under PCP not solved as a separate aggregate target in this map.
  \item (18): wage Phillips curve (\texttt{union\_wage\_nkpc} target residual).
  \item (19)--(20): monetary policy rule (\texttt{monetary\_rule\_taylor} target residual).
  \item (21): domestic goods market clearing (\texttt{goods\_market} target residual).
  \item (22): net foreign asset definition for IRF accounting.
  \item (23): current-account law of motion for IRF accounting.
  \item Additional helper block: inflation definitions.
\end{itemize}

\subsection{Quantitative Equations (47)--(53): Mapping Notes}
\begin{itemize}
  \item (47): Stone-Geary non-homothetic CES aggregator replaces base equation (3).
  \item (48): heterogeneous labor-shock incidence enters as household inputs.
  \item (49): home-goods Phillips curve (\texttt{piH\_nkpc} target residual).
  \item (50): foreign-goods Phillips curve (inactive in paper calibration when $\theta_F=0$).
  \item (51): home-goods-foreign Phillips curve (\texttt{piHstar\_nkpc} target residual).
  \item (52)--(55): delayed substitution system (\texttt{delayed\_substitution} with four targets), appendix D3.
\end{itemize}

\subsection{Additional Notes}
\begin{itemize}
  \item Portfolio choice is indeterminate in the one-asset setup.
  \item Rest of world is representative-agent with an exogenous patience shock.
  \item PCP baseline with no foreign price rigidity.
  \item No government spending, taxation, or debt in the core setup.
  \item Main solution emphasizes aggregate outcomes, with quantitative add-ons for distributional channels.
\end{itemize}

\clearpage
\section{Code Listing}
The model implementation can be read directly from the source file below.

\lstinputlisting[
  language=Python,
  caption={\texttt{arss\_hank\_quant.py}},
  label={lst:arss_hank_quant}
]{../arss_hank_quant.py}

\appendix
\section{Appendix Figures}
\begin{figure}[h!]
  \centering
  \includegraphics[width=\linewidth,height=0.75\textheight,keepaspectratio]{../figures/ha_oe_quant_dag.png}
  \caption{Quantitative model DAG}
\end{figure}

\begin{figure}[h!]
  \centering
  \includegraphics[width=\linewidth,height=0.75\textheight,keepaspectratio]{../figures/ha_oe_base_irf.png}
  \caption{Base model IRFs}
\end{figure}

\end{document}
